\documentclass{article}

\usepackage{titlesec}
\usepackage{tikz-among-us}

\author{Impostor from among us}
\usepackage{tikz}
\usetikzlibrary{positioning}

\title{\begin{tikzpicture}
			\amongUsI[scale=0.1, shift={(50.5,0)},
			rotate around={0:(1.75,2.3)}]
			{blue}{cyan}
		\end{tikzpicture}
		SASy cheatsheat
		\begin{tikzpicture}
			\amongUsI[scale=0.1, shift={(50.5,0)},xscale = -1,
			rotate around={0:(1.75,2.3)}]
			{red}{cyan}
		\end{tikzpicture}}
\date{}



\titleformat{\subparagraph}
    {\normalfont\normalsize\bfseries}{\thesubparagraph}{1em}{}
\titlespacing*{\subparagraph}{\parindent}{3.25ex plus 1ex minus .2ex}{.75ex plus .1ex}


\begin{document}
\maketitle


\section*{týden 1}
	\subparagraph*{Prevádzanie sinu a kosínusu na komplexné funkcie a naopak}
		(Převádění sinu a cosinu na komplexní funkce a naopak)
	
		\[x(t) = Ae^{\sigma t}(\cos(\omega t+ \varphi) + j\sin(\omega t+ \varphi )) 
		= A e^{\varphi j}e^{(\sigma + \omega j)t}\]

		\[A = |A|e^{\omega j}, \quad s = \sigma + \omega j\]

		\[x(t) = Ae^{st} = A \underbrace{\cos(st)}_{Re} 
		+ j A\underbrace{sin(st)}_{Im}\]
		
		\[x(t) = \frac{e^{j\omega t} - e^{-j\omega t}}{2j} = \sin(\omega t)\]
		
		\[x(t) = \frac{e^{j\omega t} + e^{-j\omega t}}{2} = \cos(\omega t)\]
	\subparagraph*{Vzorkovací vlastnost dirakovy funkce}
		\[\int^{\infty}_{-\infty} f(t)\delta(t - t_0) dt 
		= \int^{\infty}_{-\infty} f(t_0)y\delta(t-t_0) 
		= f(t_0)\int^{\infty}_{-\infty} \delta(t - t_0) = f(t_0)\]

\section*{týden 2}
	\subparagraph*{Průměrná hodnota, výkon, energie}

		Průměrná hodnota
			\[Avg\{x[n]\} = \frac{1}{N_2-N_1+1} \sum^{N_2}_{N=N_1} x[n] 
			\qquad Avg\{x[n]\} = \lim_{N\to\infty}\frac{1}{2N+1}\sum_{n=-N}^{N}x[n]\]

			\[Avg\{x(t)\} = \frac{1}{T_2 - T_1}\int_{T_2}^{T_1}x(t)dt 
			\qquad Avg\{x(t)\} = \lim_{T\to\infty}\frac{1}{2T}\int_{T}^{-T}x(t)dt\]  
		Okamžitý výkon
			\[p(t) = |x(t)|^2 \qquad p[n] = |x[n]|^n\]
		Průměrný výkon
			\[P = Avg\{|x(t)|^2\} \qquad P = Avg\{|x[n]|^2\}\]
		Energie
			\[E = \int_\jmath p(t) dt = \int_\jmath |x(t)|^2 dt \qquad E 
			= \sum_{n\in\jmath} p[n] = \sum_{n\in\jmath}|x[n]|^2\]
	\subparagraph*{Fourierova řada}
		Reálná Fourierova řada
			\[\sum_{\omega} A(\omega) \cos(\omega t + \varphi(\omega)) 
			\quad \omega = 2\pi k f\]
		Komplexní Fourierova řada
			\[x(t) = \sum_{k = -\infty}^\infty{C_k e^{j \omega_k t}}, 
			\quad \omega_k = 2 \pi fk = \frac{2 \pi}{T}k\]
		Koeficient komplexní Fourierovy řady
			\[C_k = \frac{1}{T}\int_{<T>} x(t)e^{-j\omega_k t}dt\]



\section*{týden 3}
	\subsection*{Fourierovy transformace}
		Fourierova transformace ve spojitém čase (CTFT)
			\[\mathcal{F}\{x(t)\} = X(\omega) 
			= \int_{-\infty}^{\infty} x(t)e^{-j\omega t} dt\]

		inverzní CTFT 
			\[x(t) = \frac{1}{2\pi} 
			\int_{-\infty}^{\infty} X(\omega)e^{j\omega t}d\omega\]

		Fourierova transformace v diskrétním čase (DTFT) 
		\[X(\Omega) = \sum_{n = - \infty}^{\infty} x[n]e^{-j\Omega n}\]

		inverzní (DTFT) 
		\[x[n] = \frac{1}{2 \pi} \int_{0}^{2 \pi} X(\Omega ) e^{j \Omega n} d\Omega\]


		Diskrétní Fourierova Transformace (DFT)
		\[X[k] = \sum_{s=0}^{N-1} x[n] e^{-j \frac{2\pi}{N}kn}\]

		Inverzní DFT
		\[x_p[n] = \frac{1}{N} \sum_{s=0}^{N-1} X[k] e^{j \frac{2\pi}{N}kn}\]


\section*{týden 4}
	\subsection*{Lineární konvoluce}
		\[x[n] \ast y[n] = \sum_{m=-\infty}^{\infty}x[n-m]y[m]
		= \sum_{m = -\infty}^{\infty}x[m]y[n-m]\]
		\[x(t) \ast y(t) = \int_{-\infty}^{\infty} x(t-\tau)y(\tau)d\tau
		= \int_{-\infty}^{\infty}x(\tau) y(t-\tau)d\tau\]

		Fir filtry(finite impulse responce)
		
	
\section*{týden 5}
\section*{týden 6}
\section*{týden 7}
\section*{týden 8}
\section*{týden 9}
	\subsubsection*{Z-transformace}
		Definice
		\[\mathcal{Z}\{x[n]\}=\sum_{n=0}^{\infty}x[n]z^{-n}\]
		Odvození vzorců
		\[\mathcal{Z}\{\delta[n]\}=\sum_{n=0}^{\infty}\delta[n]z^{-n}=z^0=1\]
		impulsní odezva se zjštuje, tak že najdeme Z/L transformaci
		a přechodovou odezvu (odezva na jednotkový skok) najdeme tak,
		že zase najdem Z/L transformaci ale pronásonbenou 
		zpětnou Z/L transformací jednotkovýho skoku


\end{document}